\chapter{Ekstremalna teoria grafów: twierdzenia Mantela, Turána i Ramseya}
\label{ch:tg-extremal-ramsey}

W rozdziale \ref{ch:tg-podstawy} wprowadziliśmy grafy proste,
kliki oraz zbiory niezależne. Zbadamy teraz największą liczbę
krawędzi przy zakazie ustalonej kliki, a następnie twierdzenia
gwarantujące istnienie kliki lub zbioru niezależnego.
Grafy losowe i ekspandery zostaną opracowane w osobnych modułach.

\section{Funkcja ekstremalna i grafy Turána}
\label{sec:tg-extremal-notation}

Niech $G=(V,E)$ będzie skończonym grafem prostym
nieskierowanym, $n=|V|$, $m=|E|$.
Oznaczamy sąsiedztwo wierzchołka $v$ przez $N_G(v)$,
jego stopień przez $d_G(v)=|N_G(v)|$, a liczbę krawędzi
grafu $G$ również przez $e(G)=|E|$.

\begin{definicja}[Funkcja ekstremalna]
\label{def:tg-extremal-ex}
Dla skończonego grafu prostego nieskierowanego $H$ i liczby całkowitej
$n\geq0$ niech $\operatorname{ex}(n,H)$ oznacza maksymalną
liczbę krawędzi grafu prostego o $n$ wierzchołkach,
który nie zawiera $H$ jako podgrafu (niekoniecznie
indukowanego). Jeśli nie istnieje żaden taki graf,
przyjmujemy $\operatorname{ex}(n,H)=-\infty$.
W przeciwnym przypadku maksimum istnieje, ponieważ
na ustalonym zbiorze $n$ wierzchołków jest dokładnie
$2^{\binom n2}$ grafów prostych. Na przykład dla zakazu
$K_1$ i $n\geq1$ rodzina dopuszczalnych grafów jest pusta.
\end{definicja}

\begin{definicja}[Graf Turána]
\label{def:tg-turan}
Dla liczb całkowitych $r\geq1$ i $n\geq0$ grafem Turána $T_r(n)$
nazywamy pełny graf $r$-dzielny na $n$ wierzchołkach,
którego $r$ części mają liczności różniące się
o co najwyżej jeden; dopuszczamy części puste.
Niech $n=ar+q$ będzie wynikiem dzielenia z resztą,
gdzie $a,q$ są całkowite, $a\geq0$ i $0\leq q<r$.
Ponieważ liczności części różnią się o najwyżej jeden,
ich minimum jest równe $a$: gdyby było większe, suma
przekroczyłaby $n$, a gdyby było mniejsze, suma
byłaby mniejsza od $ar$ (co najmniej jedna część
miałaby liczność mniejszą od $a$).
Zatem części mają liczności $a$ lub $a+1$.
Z warunku sumy $n=ar+q$ wynika, że dokładnie $q$
części ma liczność $a+1$, a pozostałe $r-q$ liczność $a$.
\end{definicja}

Każda klika $T_r(n)$ zawiera najwyżej jeden wierzchołek
z każdej części, więc graf nie zawiera $K_{r+1}$.
Niech $n_1,\ldots,n_r$ będą licznościami części.
Każdą krawędź liczymy dokładnie raz w sumie
$\sum_{i<j}n_in_j=\tfrac12((\sum_i n_i)^2-\sum_i n_i^2)$.
Podstawiając liczności części z definicji, otrzymujemy
\[
e(T_r(n))=\frac12\left(n^2-q(a+1)^2-(r-q)a^2\right).
\]

\section{Twierdzenie Mantela}
\label{sec:tg-extremal-mantel}

\begin{twierdzenie}[Twierdzenie Mantela]
\label{thm:tg-mantel}
Dla każdej liczby całkowitej $n\geq0$ zachodzi
\[
\operatorname{ex}(n,K_3)=\left\lfloor\frac{n^2}{4}\right\rfloor.
\]
Równość osiąga graf pełny dwudzielny
$K_{\lfloor n/2\rfloor,\lceil n/2\rceil}$.
\end{twierdzenie}
\begin{proof}
Dla $n\leq1$ żaden graf prosty nie ma krawędzi,
a obie strony równości są zerowe.
Niech $G=(V,E)$ będzie grafem prostym nieskierowanym
na $n\geq2$ wierzchołkach bez trójkątów i niech $m=|E|$.
Jeśli $uv\in E$, to $N_G(u)\cap N_G(v)=\varnothing$,
ponieważ wspólny sąsiad tworzyłby trójkąt.
Zatem $d_G(u)+d_G(v)\leq n$.
W sumie stopni końców wszystkich krawędzi składnik
$d_G(v)$ występuje dokładnie $d_G(v)$ razy, ponieważ
tyle krawędzi jest incydentnych z $v$. Stąd
\[
\sum_{v\in V}d_G(v)^2
=\sum_{uv\in E}\bigl(d_G(u)+d_G(v)\bigr)\leq nm.
\]
Lemat o uściskach dłoni i nierówność
Cauchy'ego--Schwarza dają
\[
\sum_{v\in V}d_G(v)^2
\geq\frac{(\sum_vd_G(v))^2}{n}=\frac{4m^2}{n}.
\]
Łącząc nierówności, mamy $4m^2/n\leq nm$.
Gdy $m>0$, po pomnożeniu przez $n$ i podzieleniu
przez $4m$ otrzymujemy $m\leq n^2/4$;
dla $m=0$ nierówność jest natychmiastowa.
Ponieważ $m$ jest całkowite,
$m\leq\lfloor n^2/4\rfloor$.
Graf o wskazanych dwóch częściach
nie zawiera trójkąta, a ma dokładnie
$\lfloor n/2\rfloor\lceil n/2\rceil
=\lfloor n^2/4\rfloor$ krawędzi: dla $n=2j$ oba wyrażenia
wynoszą $j^2$, a dla $n=2j+1$ wynoszą $j(j+1)$.
\end{proof}

\begin{przyklad}[Osiem wierzchołków]
\label{ex:tg-mantel-8}
Graf na ośmiu wierzchołkach bez trójkąta
może mieć co najwyżej $16$ krawędzi.
Graf $K_{4,4}$ osiąga tę liczbę.
Dołożenie krawędzi wewnątrz którejkolwiek
części tworzy trójkąt z każdym wierzchołkiem
drugiej części.
\end{przyklad}

\begin{figure}[htbp]
\centering
\begin{tikzpicture}[scale=.9]
\foreach \i/\y in {1/1.2,2/0,3/-1.2}{
  \node[tgVertex] (a\i) at (-2,\y) {$a_{\i}$};
  \node[tgVertex] (b\i) at (2,\y) {$b_{\i}$};
}
\foreach \i in {1,2,3}{
  \foreach \j in {1,2,3}{\draw[tgEdge] (a\i)--(b\j);}
}
\node[font=\small] at (-2,-1.8) {$A$};
\node[font=\small] at (2,-1.8) {$B$};
\end{tikzpicture}
\caption{Graf $K_{3,3}$ ma dziewięć krawędzi
i nie zawiera trójkąta; każda krawędź łączy obie części.}
\label{fig:tg-mantel-graph}
\end{figure}

\section{Symetryzacja Zykowa i twierdzenie Turána}
\label{sec:tg-extremal-turan}

Dwa różne wierzchołki $u,v$ nazywamy
\emph{bliźniakami}, jeśli $N_G(u)=N_G(v)$.
Jeśli byłyby sąsiednie, mielibyśmy $u\in N_G(v)=N_G(u)$,
co przeczy brakowi pętli. Równość zbiorów sąsiadów
definiuje relację równoważności także dla $u=v$;
każda jej klasa jest zbiorem niezależnym.
Pomiędzy dwiema różnymi klasami albo nie ma żadnej
krawędzi, albo występują wszystkie możliwe krawędzie:
jeśli $ab$ jest krawędzią dla $a\in A$, $b\in B$,
to każdy $a'\in A$ jest sąsiadem $b$, a następnie
każdy $b'\in B$ jest sąsiadem $a'$.

\begin{lemat}[Symetryzacja Zykowa]
\label{lem:tg-zykov}
Niech $r\geq1$ będzie liczbą całkowitą, a $G$
skończonym grafem prostym nieskierowanym bez $K_{r+1}$.
Niech $A,B$ będą dwiema różnymi klasami bliźniaków,
pomiędzy którymi nie ma krawędzi.
Zastąpmy sąsiedztwa każdego wierzchołka $B$
sąsiedztwami jednego ustalonego wierzchołka $A$,
pozostawiając $A\cup B$ jako zbiór niezależny.
Nowy graf nadal nie zawiera $K_{r+1}$,
a liczba krawędzi zmienia się o
$|B|(d_G(a)-d_G(b))$, gdzie $a\in A$, $b\in B$.
Jeśli oba stopnie są równe, operacja nie zmienia
liczby krawędzi, scala $A,B$ w jedną klasę
bliźniaków i nie rozdziela żadnej innej klasy.
\end{lemat}
\begin{proof}
Każda klika po modyfikacji zawiera
co najwyżej jeden wierzchołek $B$.
Jeśli zawiera taki wierzchołek, nie zawiera
żadnego elementu $A$, gdyż $A\cup B$ jest niezależny.
Możemy więc zastąpić ten wierzchołek dowolnym
elementem $A$. Wszystkie pozostałe krawędzie
kliki są niezmienione, a nowe sąsiedztwo $B$
jest pierwotnym sąsiedztwem elementu $A$.
Otrzymujemy klikę tej samej wielkości w grafie
pierwotnym. Klika niezawierająca elementów $B$
pozostaje bez zmian.
Nie ma krawędzi wewnątrz $B$ ani między
$A$ i $B$, więc zmiana liczby krawędzi
jest sumą zmian stopni wierzchołków $B$:
$|B|(d_G(a)-d_G(b))$.
Po symetryzacji $A\cup B$ ma identyczne
sąsiedztwa. Jeśli $x,y$ spoza $A\cup B$
miały wcześniej jednakowe sąsiedztwa,
były jednocześnie sąsiadami każdego
$a\in A$ albo żaden z nich nie był.
Zmiany relacji do $B$ są zatem dla $x,y$
identyczne i nie rozbijają ich klasy.
\end{proof}

\begin{twierdzenie}[Twierdzenie Turána]
\label{thm:tg-turan}
Dla wszystkich liczb całkowitych $n\geq0$, $r\geq1$ zachodzi
\[
\operatorname{ex}(n,K_{r+1})=e(T_r(n)).
\]
\end{twierdzenie}
\begin{proof}
Dla $n=0$ jedynym grafem jest graf pusty,
więc teza zachodzi. Niech $n>0$.
Wśród wszystkich grafów na ustalonym zbiorze
$n$ wierzchołków bez $K_{r+1}$ (jest wśród nich
graf bez krawędzi) wybierzmy graf $G$
o największej liczbie krawędzi,
a wśród takich grafów ten, dla którego
suma kwadratów liczności klas bliźniaków
jest największa. Obie maksymalizacje mają
rozwiązania, gdyż rozważamy zbiór skończony.

Jeżeli dwa niesąsiadujące wierzchołki
należałyby do różnych klas bliźniaków
$A,B$, to nie byłoby krawędzi między
tymi klasami. Gdy ich stopnie są różne,
symetryzacja klasy o mniejszym stopniu
do klasy o większym zwiększa liczbę krawędzi,
co przeczy pierwszej maksymalizacji.
Gdy stopnie są równe, lemat
\ref{lem:tg-zykov} pozwala scalić klasy
bez zmiany liczby krawędzi
i bez rozbicia pozostałych klas.
Po scaleniu $A,B$ wkład tych klas rośnie o
$(|A|+|B|)^2-|A|^2-|B|^2=2|A||B|>0$.
Żadna inna klasa się nie rozdziela, a ewentualne
dalsze scalenia nie zmniejszają sumy kwadratów.
Suma ta rośnie więc co najmniej o $2|A||B|$,
wbrew drugiej maksymalizacji.
Zatem każde dwie różne klasy bliźniaków
są połączone wszystkimi możliwymi
krawędziami, a $G$ jest pełnym grafem
wielodzielnym.

Niech $k$ oznacza liczbę niepustych części.
Wybierając po jednym wierzchołku z każdej,
dostajemy klikę $K_k$, więc $k\leq r$.
Jeśli $k<\min\{n,r\}$, to pewna część
ma co najmniej dwa wierzchołki.
Możemy ją rozdzielić na dwie niepuste
części i dodać krawędzie między nimi.
Nowy graf jest pełny $(k+1)$-dzielny, gdzie
$k+1\leq r$, więc nie zawiera $K_{r+1}$.
Jeśli nowe części mają liczności $p,q>0$,
dodaliśmy dokładnie $pq>0$ krawędzi, sprzecznie
z maksymalnością liczby krawędzi.
Stąd $k=\min\{n,r\}$.

Oznaczmy rozmiary części przez $n_1,\ldots,n_k$.
Liczba krawędzi wynosi
$\frac12(n^2-\sum_i n_i^2)$.
Jeśli $n_i\geq n_j+2$,
zastępujemy graf pełny $k$-dzielny grafem
o licznościach części $n_i-1,n_j+1$,
pozostawiając pozostałe liczności bez zmian.
Nowy graf nadal ma najwyżej $r$ części, więc
nie zawiera $K_{r+1}$. Suma kwadratów maleje o
$n_i^2+n_j^2-(n_i-1)^2-(n_j+1)^2
=2(n_i-n_j-1)>0$, a więc liczba krawędzi rośnie.
To niemożliwe, zatem części mają rozmiary
różniące się co najwyżej o jeden.
Po dołączeniu $r-k$ pustych części otrzymaliśmy
graf izomorficzny z $T_r(n)$, a zatem
maksymalna liczba krawędzi jest równa $e(T_r(n))$.
Dowód pokazuje, że pewien graf ekstremalny ma tę postać;
sam wybór grafu przez dodatkową maksymalizację
nie dowodzi jeszcze jednoznaczności wszystkich
grafów ekstremalnych.
\end{proof}

\begin{stwierdzenie}[Jednoznaczność grafu ekstremalnego w twierdzeniu Turána]
\label{prop:tg-turan-unique}
Jeśli graf prosty $G$ o $n$ wierzchołkach
nie zawiera $K_{r+1}$ i ma $e(T_r(n))$ krawędzi,
to jest izomorficzny z $T_r(n)$.
\end{stwierdzenie}
\begin{proof}
Dowód przeprowadzimy przez indukcję względem $r$.
Dla $r=1$ graf nie zawiera żadnej krawędzi.
Przypadek $n=0$ jest również bezpośredni.
Niech $r\geq2$, $n>0$ i wybierzmy wierzchołek
$v$ maksymalnego stopnia $d$. Oznaczmy
$A=N_G(v)$ i $B=V(G)\setminus A$.
Graf indukowany na $A$ nie zawiera $K_r$,
gdyż wraz z $v$ taka klika tworzyłaby $K_{r+1}$.
Z twierdzenia \ref{thm:tg-turan} mamy więc
$e(G[A])\leq e(T_{r-1}(d))$.
Ponadto
\[
e(G)=e(G[A])+\sum_{b\in B}d_G(b)-e(G[B])
\leq e(T_{r-1}(d))+|B|d-e(G[B]).
\]
Graf $H$ otrzymany z $T_{r-1}(d)$ i niezależnego
zbioru $B$ przez dodanie wszystkich krawędzi
między nimi nie zawiera $K_{r+1}$, a jego
liczba krawędzi wynosi $e(T_{r-1}(d))+|B|d$.
Stąd $e(H)\leq e(T_r(n))=e(G)$.
W połączeniu z poprzednią nierównością otrzymujemy
$e(G[B])=0$ oraz równość we wszystkich oszacowaniach.
W szczególności każdy $b\in B$ ma stopień $d$,
a ponieważ $B$ jest niezależny i $|A|=d$,
jest sąsiadem każdego wierzchołka $A$.
Z hipotezy indukcyjnej $G[A]\cong T_{r-1}(d)$.
Graf $G$ jest zatem pełnym grafem $r$-dzielnym,
z dopuszczeniem pustych części.
Jeśli dwie jego części mają liczności różniące
się co najmniej o dwa, przeniesienie jednego
wierzchołka z większej do mniejszej w konstrukcji
grafu pełnego wielodzielnego zwiększa liczbę
krawędzi, jak w dowodzie twierdzenia
\ref{thm:tg-turan}. Przeczy to ekstremalności.
Wszystkie liczności różnią się więc co najwyżej
o jeden i $G\cong T_r(n)$.
\end{proof}

\begin{wniosek}[Prostsze oszacowanie]
\label{cor:tg-turan-bound}
Niech $n\geq0$, $r\geq1$ będą liczbami całkowitymi,
a $G$ grafem prostym nieskierowanym na $n$
wierzchołkach bez $K_{r+1}$. Wówczas
\[
|E(G)|\leq\frac12\left(1-\frac1r\right)n^2.
\]
\end{wniosek}
\begin{proof}
Jeżeli liczności $r$ części grafu Turána
(z uwzględnieniem pustych) wynoszą $n_i$,
to nierówność Cauchy'ego--Schwarza daje
$\sum_i n_i^2\geq n^2/r$.
Na mocy twierdzenia \ref{thm:tg-turan} oraz wzoru
na liczbę krawędzi grafu wielodzielnego otrzymujemy
\[
|E(G)|\leq e(T_r(n))
=\tfrac12(n^2-\textstyle\sum_i n_i^2)
\leq\tfrac12(1-1/r)n^2.
\]
\end{proof}

\begin{przyklad}[Zakaz $K_4$]
\label{ex:tg-turan-ten}
Dla $n=10$ i $r=3$ podział grafu
$T_3(10)$ ma części wielkości $4,3,3$.
Ma on $\tfrac12(100-16-9-9)=33$ krawędzie.
Żaden dziesięciowierzchołkowy graf bez $K_4$
nie ma więcej krawędzi.
\end{przyklad}

\section{Nierówność Caro--Weiego i dolne oszacowania niezależności}
\label{sec:tg-caro-wei}

\begin{twierdzenie}[Nierówność Caro--Weiego]
\label{thm:tg-caro-wei}
Niech $G=(V,E)$ będzie skończonym
grafem prostym nieskierowanym o $n\geq1$
wierzchołkach i $m$ krawędziach.
Przez $\alpha(G)$ oznaczmy maksymalną
liczność zbioru niezależnego w $G$.
Wówczas
\[
\alpha(G)\geq
\sum_{v\in V}\frac{1}{d_G(v)+1}
\geq\frac{n^2}{n+2m}.
\]
Oba oszacowania są ostre; równość
w pierwszej nierówności zachodzi
między innymi dla sum rozłącznych
grafów pełnych.
\end{twierdzenie}
\begin{proof}
Wybierzmy jednostajnie losowo
permutację wszystkich $n$ wierzchołków.
Do zbioru $I$ włączmy wierzchołek $v$
wtedy i tylko wtedy, gdy $v$
występuje wcześniej niż wszyscy
jego sąsiedzi w tej permutacji.
Zbiór $I$ jest niezależny:
gdyby dwa sąsiednie wierzchołki
$u,v$ należały do $I$,
pierwszy z nich musiałby
występować wcześniej od
drugiego i jednocześnie
drugi wcześniej od pierwszego,
co jest niemożliwe.

Niech $S_v=\{v\}\cup N_G(v)$ oraz $h=|S_v|=d_G(v)+1$.
Każdą ustaloną kolejność elementów $S_v$
można rozszerzyć do dokładnie
$\binom nh(n-h)!=n!/h!$ permutacji zbioru $V$:
wybieramy $h$ pozycji dla $S_v$ i dowolnie
porządkujemy pozostałe elementy.
Względna kolejność elementów $S_v$ jest więc
jednostajnie losowa.
Wierzchołek $v$ jest
najwcześniejszy w tym
$(d_G(v)+1)$-elementowym
zbiorze z prawdopodobieństwem
$1/(d_G(v)+1)$.
Dla losowego zbioru $I$ mamy
$|I|=\sum_{v\in V}\chi_I(v)$.
Ponieważ $\mathbb E\chi_I(v)=\mathbb P(v\in I)$,
z liniowości wartości oczekiwanej, bez potrzeby
niezależności zdarzeń dla różnych wierzchołków,
\[
\mathbb E|I|
=\sum_{v\in V}\frac{1}{d_G(v)+1}.
\]
Maksimum liczności $I$ po skończonym zbiorze
permutacji jest nie mniejsze niż średnia tych
liczności. Istnieje zatem permutacja, dla której
$|I|\geq\mathbb E|I|$.
Ponieważ każdy taki $I$
jest niezależny, otrzymujemy
pierwszą nierówność.

Z nierówności Cauchy'ego--Schwarza
oraz lematu o uściskach dłoni,
\[
n^2=
\left(\sum_{v\in V}
\sqrt{d_G(v)+1}\,
\frac1{\sqrt{d_G(v)+1}}\right)^2
\leq(n+2m)
\sum_{v\in V}\frac1{d_G(v)+1}.
\]
Dzieląc przez $n+2m>0$,
dostajemy drugą nierówność.

Jeżeli $G$ jest sumą rozłączną
grafów pełnych $K_{a_1},\ldots,K_{a_k}$
o dodatnich licznościach $a_i$,
to każdy zbiór niezależny zawiera
co najwyżej jeden wierzchołek
z każdej składowej. Wybór dokładnie jednego
wierzchołka z każdej składowej daje zbiór
niezależny, ponieważ pomiędzy składowymi
nie ma krawędzi. Stąd $\alpha(G)=k$. Z drugiej strony
suma z tezy wynosi
$\sum_i a_i/a_i=k$.
Dla grafu pełnego $K_n$ także
druga nierówność jest równością,
bo $m=\binom n2$ i
$n^2/(n+2m)=1=\alpha(K_n)$.
\end{proof}

\begin{wniosek}[Stopień średni a zbiór niezależny]
\label{cor:tg-caro-wei-average}
Niech $G=(V,E)$ będzie skończonym grafem prostym
nieskierowanym, $n=|V|\geq1$, $m=|E|$ oraz
$\alpha(G)$ liczbą niezależności.
Dla średniego stopnia $\overline d=2m/n$ zachodzi
\[
\alpha(G)\geq
\left\lceil\frac{n}{\overline d+1}\right\rceil.
\]
W szczególności jeśli
$\overline d\leq d$ dla pewnej
liczby rzeczywistej $d\geq0$,
to $\alpha(G)\geq\lceil n/(d+1)\rceil$.
\end{wniosek}
\begin{proof}
Z Twierdzenia~\ref{thm:tg-caro-wei}
otrzymujemy
$\alpha(G)\geq n/(\overline d+1)$.
Ponieważ $\alpha(G)$ jest liczbą
całkowitą, możemy przejść do
sufitu. Drugą tezę daje
monotoniczność mianownika.
\end{proof}

\begin{przyklad}[Suma dwóch klik]
\label{ex:tg-caro-wei-two-cliques}
Niech $G=K_3\sqcup K_5$.
Wówczas $n=8$, $m=3+10=13$
i $\alpha(G)=2$.
Suma Caro--Weiego wynosi
$3/3+5/5=2$, a
oszacowanie zależne
tylko od liczby krawędzi
daje $n^2/(n+2m)=64/34<2$.
Pierwsza nierówność jest w tym przykładzie
równością, a druga jest ścisła. Po zaokrągleniu
w górę obie dają jednak tę samą całkowitą
dolną granicę $2$.
\end{przyklad}

\section{Liczby Ramseya i skończone twierdzenie Ramseya}
\label{sec:tg-extremal-ramsey}

\begin{definicja}[Liczba Ramseya]
\label{def:tg-ramsey-number}
Dla liczb całkowitych $s,t\geq1$
liczbą Ramseya $R(s,t)$ nazywamy
najmniejszą dodatnią liczbę całkowitą $N$,
taką że każde dwukolorowanie krawędzi
$K_N$ na czerwono i niebiesko zawiera
czerwoną klikę $K_s$ albo niebieską
klikę $K_t$. Klikę jednoelementową
traktujemy jako jednobarwną.
Istnienie takiego $N$ udowodnimy poniżej.
Jeżeli własność zachodzi dla $N$, zachodzi też
dla każdej liczby $n\geq N$, ponieważ dowolne
kolorowanie $K_n$ można ograniczyć do $N$ wierzchołków.
\end{definicja}

\begin{twierdzenie}[Skończone twierdzenie Ramseya]
\label{thm:tg-ramsey-finite}
Dla wszystkich liczb całkowitych $s,t\geq1$
liczba $R(s,t)$ istnieje i jest skończona.
Zachodzą równości
$R(1,t)=R(s,1)=1$, a dla $s,t\geq2$
nierówność
\[
R(s,t)\leq R(s-1,t)+R(s,t-1).
\]
\end{twierdzenie}
\begin{proof}
Indukujemy po $s+t$.
Jeśli $s=1$ albo $t=1$, to dowolne kolorowanie
$K_1$ zawiera wymaganą klikę jednoelementową;
najmniejsza dodatnia liczba to $1$.
Niech $s,t\geq2$ i
$N=R(s-1,t)+R(s,t-1)$.
W dowolnym dwukolorowaniu $K_N$
wybierzmy wierzchołek $v$.
Ma on $N-1$ sąsiadów,
więc co najmniej
$R(s-1,t)$ spośród krawędzi
wychodzących z $v$ jest
czerwonych albo co najmniej
$R(s,t-1)$ jest niebieskich:
w przeciwnym razie ich suma
nie przekroczyłaby $N-2$.

Jeżeli pierwsza alternatywa zachodzi, wybieramy
$R(s-1,t)$ czerwonych sąsiadów $v$.
Z założenia indukcyjnego liczba $R(s-1,t)$
istnieje, więc na mocy jej definicji wśród
wybranych wierzchołków istnieje czerwona
$K_{s-1}$ lub niebieska $K_t$.
Pierwszą rozszerzamy
o wierzchołek $v$
do czerwonej $K_s$.
Jeśli druga alternatywa zachodzi, wybieramy
$R(s,t-1)$ niebieskich sąsiadów $v$.
Z założenia indukcyjnego i definicji $R(s,t-1)$
wśród nich istnieje czerwona $K_s$ albo
niebieska $K_{t-1}$;
drugą rozszerzamy o $v$.
Tym samym $N$ ma
wymaganą własność,
co dowodzi skończoności
i nierówności rekurencyjnej.
\end{proof}

\begin{wniosek}[Dwumianowe oszacowanie Ramseya]
\label{cor:tg-ramsey-binomial}
Dla wszystkich liczb całkowitych $s,t\geq1$ zachodzi
\[
R(s,t)\leq\binom{s+t-2}{s-1}.
\]
W szczególności
$R(k,k)\leq\binom{2k-2}{k-1}$.
\end{wniosek}
\begin{proof}
Jeżeli $s=1$ lub $t=1$,
obie strony są równe $1$.
Dla $s,t\geq2$ stosujemy indukcję względem sumy
$s+t$. Nierówność z twierdzenia
\ref{thm:tg-ramsey-finite}, hipoteza indukcyjna
dla par $(s-1,t)$ i $(s,t-1)$ oraz tożsamość
Pascala dają
\[
\begin{aligned}
R(s,t)&\leq R(s-1,t)+R(s,t-1)\\
&\leq\binom{s+t-3}{s-2}+\binom{s+t-3}{s-1}\\
&=\binom{s+t-2}{s-1}.
\end{aligned}
\]
\end{proof}

\begin{stwierdzenie}[Wersja dla kliki i zbioru niezależnego]
\label{prop:tg-ramsey-clique-independent}
Niech $s,t\geq1$ i $n\geq R(s,t)$ będą liczbami
całkowitymi. Każdy graf prosty nieskierowany
o $n$ wierzchołkach
zawiera klikę o $s$ wierzchołkach
albo zbiór niezależny
o $t$ wierzchołkach.
\end{stwierdzenie}
\begin{proof}
W grafie pełnym $K_n$ kolorujemy
na czerwono krawędzie
danego grafu $G$,
a na niebiesko pozostałe.
Klika czerwona jest
kliką $G$, a klika
niebieska jest
zbiorem niezależnym w $G$.
Stosujemy definicję $R(s,t)$.
\end{proof}

\section{Dokładna wartość \texorpdfstring{$R(3,3)$}{R(3,3)}}
\label{sec:tg-extremal-r33}

\begin{twierdzenie}[Liczba Ramseya dla trójkątów]
\label{thm:tg-ramsey-33}
Zachodzi $R(3,3)=6$.
Oznacza to, że każdy
graf prosty o sześciu
wierzchołkach zawiera
trójkąt albo trójelementowy
zbiór niezależny,
podczas gdy dla pięciu
wierzchołków obie
konfiguracje można wykluczyć.
\end{twierdzenie}
\begin{proof}
Najpierw zauważmy, że
$R(2,3)=3$: w dwukolorowaniu
$K_3$ albo istnieje
czerwona krawędź,
albo wszystkie krawędzie
są niebieskie, a
w $K_2$ niebieska
kolorystyka wyklucza
oba żądane przypadki.
Analogicznie $R(3,2)=3$.
Zatem
$R(3,3)\leq R(2,3)+R(3,2)=6$.

Na zbiorze wierzchołków $\{0,1,2,3,4\}$
pokolorujmy krawędzie cyklu $C_5$ na czerwono,
a pozostałe krawędzie
$K_5$ na niebiesko.
Graf czerwony nie ma
trójkątów, a jego
dopełnienie także jest cyklem pięciowierzchołkowym:
jego krawędzie łączą $i$ z $i+2$ modulo $5$,
co daje cykl $0,2,4,1,3,0$. Dlatego graf niebieski
również nie ma trójkątów.
Stąd $R(3,3)>5$.
Połączenie obu
nierówności kończy dowód.
\end{proof}

\begin{figure}[htbp]
\centering
\begin{tikzpicture}[scale=.95]
\foreach \i in {0,...,4}{
  \pgfmathsetmacro{\angle}{90+72*\i}
  \node[tgVertex] (v\i) at (\angle:1.8) {$\i$};
}
\foreach \i in {0,...,4}{
  \pgfmathtruncatemacro{\j}{mod(\i+1,5)}
  \pgfmathtruncatemacro{\k}{mod(\i+2,5)}
  \draw[tgTree] (v\i)--(v\j);
  \draw[tgGhost,densely dashed] (v\i)--(v\k);
}
\end{tikzpicture}
\caption{Dwukolorowanie $K_5$ bez
jednobarwnego trójkąta.
Krawędzie ciągłe tworzą $C_5$,
a przerywane jego dopełnienie,
również izomorficzne z $C_5$.}
\label{fig:tg-ramsey-c5}
\end{figure}

\begin{przyklad}[Dwa różne zagadnienia]
\label{ex:tg-ramsey-versus-turan}
Twierdzenie Mantela mówi,
że graf o sześciu wierzchołkach
bez trójkąta ma co najwyżej
$\lfloor36/4\rfloor=9$ krawędzi.
Twierdzenie Ramseya daje
inny rodzaj wniosku:
każdy taki graf ma
trójelementowy zbiór
niezależny. Samo
oszacowanie liczby
krawędzi nie jest
jeszcze tym twierdzeniem.
\end{przyklad}

\section{Twierdzenie Ramseya dla wielu kolorów}
\label{sec:tg-extremal-multicolor}

\begin{twierdzenie}[Wielobarwne twierdzenie Ramseya]
\label{thm:tg-ramsey-multicolor}
Niech $q\geq2$ i
$k_1,\ldots,k_q\geq1$
będą liczbami całkowitymi.
Istnieje liczba $N$ taka,
że każde kolorowanie
krawędzi $K_N$ jednym
z $q$ kolorów zawiera,
dla pewnego $i\in\{1,\ldots,q\}$,
jednobarwną klikę
$K_{k_i}$ w kolorze $i$.
\end{twierdzenie}
\begin{proof}
Indukujemy po $q$.
Dla $q=2$ teza
jest twierdzeniem
\ref{thm:tg-ramsey-finite}.
Niech $q\geq3$ i załóżmy prawdziwość tezy
dla $q-1$ kolorów oraz dowolnych dodatnich
całkowitych rozmiarów klik.
Zastosujmy ją do $(k_1,\ldots,k_{q-1})$,
otrzymując dodatnią liczbę całkowitą $M$.
Bierzemy $N=R(M,k_q)$.
W zadanym kolorowaniu
$K_N$ łączymy
kolory $1,\ldots,q-1$
w jeden kolor czerwony,
a kolor $q$
zastępujemy niebieskim.
W wersji dwukolorowej
istnieje niebieska
$K_{k_q}$, która
kończy dowód, albo
czerwona $K_M$.
W drugim przypadku ograniczamy oryginalne
kolorowanie do $M$ wierzchołków czerwonej kliki.
Każda krawędź ma jeden z pierwszych $q-1$
oryginalnych kolorów, więc z wyboru $M$
istnieje $i\in\{1,\ldots,q-1\}$ oraz klika
$K_{k_i}$ w oryginalnym kolorze $i$.
\end{proof}

\section{Dalsze problemy ekstremalne}
\label{sec:tg-extremal-next}

Twierdzenie Turána
dotyczy zakazu kliki
$K_{r+1}$.
Dla ogólnego podgrafu
zabronionego $H$
funkcja
$\operatorname{ex}(n,H)$
nie musi być opisana
przez grafy wielodzielne.
Nie należy też utożsamiać
problemu zakazu podgrafu
z zakazem podgrafu
\emph{indukowanego}.

Ramseyowskie oszacowanie
dwumianowe ma charakter
ogólny, ale nie
zawsze jest ostre.
Dla $R(3,3)=6$ oszacowanie dwumianowe jest
akurat dokładne, gdyż $\binom42=6$. Aby wykazać
ostrość lub nieostrość w innym przypadku, potrzebny
jest dodatkowy argument.

W kolejnym osobnym
module rozpatrzymy
model $G(n,p)$,
metody probabilistyczne
i ich zastosowanie
do dolnych oszacowań
liczb Ramseya.
Później przejdziemy
do ekspanderów,
ich własności
kombinatorycznych
i spektralnych.
